\section{Independent approximation parameters}

\begin{theorem}[A filter exponent in both norm regimes]
    \label{thm:independent_truncation_filter_exponent}
    \lean{Real.independent_truncation_filter_exponent}
    \leanok
    Let $\ell,a\ge0$, $a\le t$, $7\le\ell g$, and
    $7ta^2\le g\ell^4$. For every $0<R\le3t+4\ell$,
    one has $a\le\ell^2\sqrt{g/R}$.
\end{theorem}
\begin{proof}
    \leanok
    If $t\le\ell$, then $a\le\ell$ and $R\le7\ell$, whence
    $a^2R\le7\ell^3\le g\ell^4$.
    If $\ell\le t$, then $R\le7t$ and the fourth-power hypothesis gives
    the same inequality. Divide by $R$ and take square roots.
\end{proof}

\begin{theorem}[An integer window at an independent threshold]
    \label{thm:independent_truncation_filter_scale}
    \lean{Real.exists_independent_truncation_filter_scale}
    \leanok
    Let $0<g\le1$ and $0\le a\le t$.
    There is an integer $\ell\ge1$ such that
    \begin{align}
        \ell\le8/g+(7ta^2/g)^{1/4},
        \label{eq:independent_truncation_integer_scale}
    \end{align}
    and $a\le\ell^2\sqrt{g/R}$ for every $0<R\le3t+4\ell$.
\end{theorem}
\begin{proof}
    \leanok
    \uses{thm:independent_truncation_filter_exponent}
    Put $b=(7ta^2/g)^{1/4}$ and $\ell=\lceil7/g+b\rceil$.
    Then $7\le\ell g$, $7ta^2\le g\ell^4$ and
    $\ell\le7/g+b+1\le8/g+b$.
    Apply the two-regime exponent bound.
\end{proof}

\begin{theorem}[The quarter-power gap and accuracy scales]
    \label{thm:quarter_power_truncation_filter_scale}
    \lean{Real.exists_quarter_power_truncation_filter_scale}
    \leanok
    Let $C\ge1$, $0<g\le1$, and $a\ge0$.
    There is an integer $\ell\ge1$ such that
    \begin{align}
        \ell\le8/g+(7C)^{1/4}g^{-1/4}a^{3/4},
        \label{eq:independent_truncation_quarter_scale}
    \end{align}
    and $a\le\ell^2\sqrt{g/R}$ for every $0<R\le3Ca+4\ell$.
    This is the independent parameter choice used for the approximation
    estimate of \cite[Theorem~4.9]{Arad2016RigorousRG}; the printed
    prescription is corrected as in
    \cite{gap:alvv16_approximation_parameters}.
\end{theorem}
\begin{proof}
    \leanok
    \uses{thm:independent_truncation_filter_scale}
    Take $t=Ca$ and use
    $(7Ca^3/g)^{1/4}=(7C)^{1/4}g^{-1/4}a^{3/4}$.
\end{proof}

\begin{theorem}[One logarithmic threshold for every projection accuracy]
    \label{thm:independent_truncation_threshold_constant}
    \lean{Real.exists_independent_truncation_threshold_constant}
    \leanok
    There is a universal constant $C\ge21$ such that, for every
    $0<g\le1$ and $0<\tau\le1$, the number
    $u=1+\log(1/g)+\log(1/\tau)$ satisfies $Cu\ge21$ and
    \begin{align}
        288e^{11/8}e^{-(Cu-g)/8}\le\tau g/10^6.
        \label{eq:independent_truncation_uniform_tail}
    \end{align}
\end{theorem}
\begin{proof}
    \leanok
    Set $K=288e^{11/8}10^6$ and $C=32+8\log K$.
    Since $K\ge1$ and $u\ge1$, one has
    $(Cu-g)/8\ge\log(K/(\tau g))$.
    Exponentiating and multiplying by $288e^{11/8}$ proves the estimate.
\end{proof}
